There are going to be constructions were we would like to be able to, once the transformer prints a token, repeat it indefinitely.
We can achieve this easily, making use of some of the tools that we had defined.
First, we need to normalize the transformer. {\color{orange} i know i define normalization later on, but i think this is the place in the text where i should write this primitive.}
Second, we need to flag the token with the $\TE$.
Third, we will apply a constant function to the flagged vector.
Since we are trying to replicate a token that was printed, i.e. not part of the input, its vector will be the last one.
This is because once the transformer prints it for the first time, we will start forcing it to repeat it.
It's important that the constant function is applied in the last layers so that the vector that goes into the $\OUTPUT$ layer (which is just a projection) is the one making all the probability going to the desired token.

Notice that this only works if the token does not appear in the input word.
That case can be solved using the non-uniformity of the model.
If the input word is of size $n$, with the $\PE$ we can flag the first $n$ vectors.
Then, we can use that flag to make them $0$ in the coordinate used by the $\TE$ for flagging the token to repeat.
We make them $0$ using a constant function as described in a previous section.

\medskip

Notice that this primitive makes the assumption that a transformer uses all of its budget before halting not a problem.
This is because if it printed a stopping symbol, but it still had budget, we can force it to repeat it until the end.